% !TeX root = ../main.tex
% 章边界：固定初始扰动和供需条件，分析损伤如何传导为服务损失。
% 下列标题为待研究内容，不代表命题已经证明或实验已经完成。
\chapter{地震扰动向应急资源服务退化的传导机制}
\label{chap:degradation}

\section{固定扰动条件下的应急资源服务模型}
\subsection{地震扰动与道路状态变化的情景表征}
\subsection{固定供需条件下的可行服务集}

\ThesisBlank

\section{道路损伤与服务可达性的约束传递}
\subsection{通行时间变化与时限可达性}
\subsection{容量受限与可交付资源量}

\ThesisBlank

\section{共享容量竞争与服务缺口形成}
\subsection{多需求点之间的道路容量竞争}
\subsection{共享库存约束与服务损失分配}

\ThesisBlank

\section{供需空间分布与损伤位置的交互影响}
\subsection{资源集聚与需求分散条件下的损伤响应}
\subsection{替代通道与关键供需联系}

\ThesisBlank

\section{服务退化的性质与临界条件分析}
\subsection{可行服务集的变化与服务损失边界}
\subsection{瓶颈约束激活与损失放大条件}

\ThesisBlank

\section{服务退化传导的机制实验与泸定情景分析}
\subsection{等强度不同位置扰动的受控比较}
\subsection{道路容量与供需分布的交互检验}
\subsection{约束放松下的损失归因与关键点识别}
\subsection{不同灾损情景下的传导差异与适用边界}

\ThesisBlank

\section{本章小结}
\ThesisBlank
